%%==================== AI 工具使用声明（2026 规定第 3 条：参考文献之前）====================
%% 标题按推荐样式：黑体小四、不加粗、居左、不参与章节编号
\par\vspace{0.6em}{\noindent\heiti\zihao{-4} AI工具使用声明\par}\vspace{0.25em}

本参赛队在竞赛过程中使用了 AI 工具，主要用于代码生成与调试、语言润色与图件脚本调整，
详细使用情况见支撑材料。以下补充说明使用环节与人工核验方式，
所有 AI 输出均经队员逐处复核后方可进入论文或程序，未被采纳的内容不计入最终成果。

\begin{itemize}\setlength{\itemsep}{2pt}
  \item \textbf{使用工具：}DSH 开发环境内嵌的 GLM-5.3 大语言模型（智谱 AI，2026-04），
        以对话方式调用；未使用任何专门针对本题的题解检索或自动答题服务。
  \item \textbf{使用环节一（代码生成与调试）：}按队员给定的模型、算法与接口要求生成
        程序初稿，用于报错定位与自检脚本编写；全部数学模型的建立、变量与约束的设定、
        参数标定方案与算法流程均由队员确定，AI 生成的代码经队员逐行审查、按模型设定
        修改并运行验证后方可使用。
  \item \textbf{使用环节二（语言润色与排版检查）：}用于摘要与正文语句的通顺性修改、
        \LaTeX{} 排版问题排查与图注表述调整。技术性陈述、数据与结论由队员撰写并逐句确认。
  \item \textbf{使用环节三（图件脚本辅助）：}用于绘图脚本的样式统一与批量出图调整；
        图形中的数据全部由本队结果文件（如 \texttt{v113\_best.json}）直接读取生成，
        未使用任何外部或生成式数据。
  \item \textbf{人工核验方式：}所有数值结果均由本队三套校验程序交叉核验——严格口径
        调度复核（\texttt{dispatch\_compliant}：截止期优先与无人机、电池双资源就绪
        队列）、完整约束评估（\texttt{eval\_full}：硬时限、加权迟到与投递区等箱属区
        的区一致零违例）、电池链独立复核（\texttt{check\_battery}：逐电池串验证，
        违规为零）；另以独立审计脚本重跑确认 80/80 货箱无遗漏、完工与能耗逐位一致，
        论文数字与结果文件经多轮全面核对。AI 给出的每处建议均在程序运行与论文复核中
        验证后才采用。
  \item \textbf{责任声明：}本文模型框架、求解算法、参数标定、数据处理、结果分析与
        全部结论均由参赛队独立完成，AI 工具未参与上述实质性判断；使用情况的逐条
        记录见支撑材料《AI 工具使用详情》。
\end{itemize}